% =============================================================================
% Introduzione: il problema e la tesi
% =============================================================================
\chapter{Introduzione}

\chapterquote{Una citazione che apre il problema.}{Autore, \textit{Opera}}

\lettrine{O}{gni testo argomentativo} parte da un problema concreto e prende
posizione. Qui si presenta il contesto, si spiega perché la questione conta e
si enuncia la \textbf{tesi} in modo chiaro e verificabile.\index{tesi}

\section{Il problema}
Descrivi la situazione e la domanda a cui il testo risponde.\footnote{Le note
a piè di pagina servono per fonti e precisazioni che non interrompono il
ragionamento.}

\section{La tesi}
\begin{concetto}[Tesi]
Una frase che il resto del testo dovrà dimostrare.
\end{concetto}

\section{Il percorso}
\begin{itemize}
	\item Primo argomento a sostegno.
	\item Secondo argomento a sostegno.
	\item Obiezioni principali e risposta.
\end{itemize}
